\section{Conclusion}\label{sec:conclusion}
We asked whether the performance of graph neural networks for EEG-based dementia classification comes
from how the signal is represented, how the brain graph is built, or the model. On a public
88-subject benchmark, with the model, data and evaluation protocol held fixed and seven
representations crossed with three graph constructions over 25 paired subject-level splits, the
answer is mostly the representation: wavelet and spectral node features lead, graph construction
matters little, and the GCN does not outperform non-graph classifiers on the same features. The best
honest three-class macro-F1 was \bestFOne{}, clearly above chance ($p=\permPThreespatialP$) but far
below published figures --- a gap that window-level splitting alone closes on our own pipeline,
inflating subject-level macro-F1 to \leakPThreespatialLeakySubj{}. A blood--brain-barrier-associated
EEG marker, the paroxysmal slow-wave event, was elevated in AD and FTD under its original definition
but was explained by background slowing, and defining it relative to each person's background
reversed the effect. We release the code, seeds, frozen folds and every raw run so that each number
in this paper can be regenerated from the data.
